El backend implementa varias medidas de seguridad complementarias,
dirigidas a proteger los datos personales en reposo y en tránsito y a
registrar las operaciones realizadas.

\begin{itemize}

  \item \textbf{Cifrado de datos personales en reposo.} El \acs{dni} de
  cada usuario se cifra con \acs{aes}-256-\acs{gcm} antes de almacenarse
  en la base de datos, con \textit{nonce} único por registro. La clave
  maestra se inyecta como variable de entorno y nunca se incorpora al
  repositorio.

  \item \textbf{\textit{Hash} de contraseñas.} Las contraseñas de usuario
  se almacenan como \textit{hash} con sal mediante Argon2; la contraseña
  en texto plano se descarta inmediatamente tras el cálculo.

  \item \textbf{Autenticación y autorización.} La autenticación se
  delega a un \dfn{plugin} intercambiable (\cref{SS:DES-PLUGINS}), con
  \acs{jwt} como implementación activa. La autorización se aplica al
  contexto \dfn{multitenant} y según los permisos del rol y la
  pertenencia del usuario a la asignatura.

  \item \textbf{Protección anti-descarga de los recursos servidos a
  alumnos.} El sistema combina cuatro mecanismos como defensa en
  profundidad: \acsp{uri} temporales con expiración corta (anotaciones),
  marca de agua dinámica con datos identificativos del alumno y marca
  temporal (páginas de examen), cifrado en tránsito con clave por sesión
  (examen), y cabeceras \acs{http} anti-caché.

  \item \textbf{Log de auditoría inmutable.} Toda operación de negocio
  con impacto en datos genera un registro en una tabla
  \textit{append-only} que captura quién (actor, \acs{ip}), qué (tipo
  de evento), sobre qué entidad y cuándo. El \textit{manager} del modelo
  bloquea explícitamente las operaciones de actualización y borrado,
  garantizando la inmutabilidad del log.

  \item \textbf{Retención y borrado periódico.} Los exámenes
  pertenecientes a cursos archivados se eliminan tras un periodo
  configurable por organización mediante una tarea periódica de Celery,
  alineada con la política de mantenimiento operativo.

\end{itemize}
